% Clause 3.1 continuation: source PDF pages 23--30 / printed pages 11--18.
% Faithfully transcribed from the user-supplied IEEE Std 1516-2010 reference PDF.

\begingroup
\setlength{\parskip}{12pt}

\newpage

\defentry{Management Object Model (MOM) and Initialization Module (MIM)}{A subset of the federation object model (FOM) that contains object model template (OMT) tables that describe the High Level Architecture (HLA) MOM. The MIM also contains additional predefined HLA constructs such as object and interaction roots, datatypes, transportation types, and dimensions. HLA specifies a standard MIM that is incorporated into all FOM Document Data (FDD) automatically by the runtime infrastructure (RTI). The standard MIM can be replaced with a user-supplied MIM containing the standard MIM plus extensions.}

\defentry{message}{A change of object instance attribute value(s), an interaction, or a deletion of an existing object instance, often associated with a particular point on the High Level Architecture (HLA) time axis, as denoted by the associated timestamp.}

\defentry{null designator}{A designator reserved to indicate an empty value that is distinguishable from all other designators. A null designator is guaranteed to compare unequal to any other designator with a nonempty value.}

\defentry{object class}{A fundamental element of a conceptual representation for a federate that reflects the real world at levels of abstraction and resolution appropriate for federate interoperability. An object class is a template for a set of characteristics that is common to a group of object instances. These characteristics correspond to the class attributes that individual federates may publish and to which other federates may subscribe.}

\defentry{object instance}{A unique instantiation of an object class that is independent of all other instances of that object class. At any point during a federation execution, the state of a High Level Architecture (HLA) object instance is defined as the collection of the values of all its instance attributes.}

\defentry{object model}{A system specification defined primarily by class characteristics and relationships. The High Level Architecture (HLA) idea of an object model is similar in many ways, but not identical, to the common idea of an object model in object-oriented literature.}

\defentry{object model framework}{The rules and terminology used to describe High Level Architecture (HLA) object models.}

\defentry{order type}{A runtime infrastructure (RTI) means of ordering messages originating from multiple joined\linebreak federates that are delivered to a single joined federate. Different categories of service are defined with\linebreak different characteristics regarding whether and how an RTI orders messages that are to be delivered to a\linebreak joined federate.}

\defentry{out of scope}{Of or pertaining to an instance attribute of an object instance for which one or more of the following are not true:}
\begin{hlalist}
    \item The object instance is known to the joined federate.
    \item The instance attribute is owned by another joined federate.
    \item The instance attribute's corresponding class attribute is one of the following:
    \begin{enumerate}[label=\arabic*),leftmargin=2.2em,labelsep=0.55em,itemsep=1pt,topsep=1pt,parsep=0pt,partopsep=0pt]
  \item A subscribed attribute of the known class of the object instance, or
  \item A subscribed attribute of the known class of the object instance with regions, and the update region set of the instance attribute at the owning joined federate overlaps the subscription region set of the instance attribute's corresponding class attribute at the known class of the instance attribute at the subscribing joined federate.
    \end{enumerate}
\end{hlalist}
\vspace{-12pt}
\hlaNote{See 6.1 of IEEE Std 1516.1-2010.}

\defentry{overlap}{}
\begin{hlalist}
    \item Pertaining to region sets: given two region sets, to have a region in each set such that the two regions overlap.
\end{hlalist}

\newpage

\begin{hlalist}
    \setcounter{enumi}{1}
    \item Pertaining to regions: given two regions, to have at least one dimension in common if and only if, for each dimension the regions have in common, their respective ranges in the common dimension overlap. If two regions do not have any dimensions in common, the regions do not overlap.
    \item Pertaining to ranges: given two ranges ($A = [a_{\mathrm{lower}}, a_{\mathrm{upper}})$ and $B = [b_{\mathrm{lower}}, b_{\mathrm{upper}})$), to overlap if and only if either $a_{\mathrm{lower}} = b_{\mathrm{lower}}$ or ($a_{\mathrm{lower}} < b_{\mathrm{upper}}$ and $b_{\mathrm{lower}} < a_{\mathrm{upper}}$).
\end{hlalist}
\hlaNote{See 9.1 of IEEE Std 1516.1-2010.}
\vspace{-3pt}

\defentry{owned}{Pertaining to the relationship between an instance attribute and a joined federate: when the joined federate has the unique right to update the instance attribute's value.}

\defentry{owned instance attribute}{An instance attribute that is explicitly modeled by the owning joined federate. A joined federate that owns an instance attribute has the unique responsibility to provide values for that instance attribute to the federation, through the runtime infrastructure (RTI), as documented in the federation object model (FOM) Document Data (FDD).}

\defentry{pair}{A grouping of two related elements (a first component and a second component), the combination of which is treated as an entity. An example of a pair would be an instance attribute grouped with its current value.}

\defentry{parameter}{A named characteristic of an interaction.}

\defentry{passel}{A group of attribute handle/value pairs from an \hlaService{Update Attribute Values} service invocation that are delivered together via a \hlaService{Reflect Attribute Values \textdagger}\ service invocation. All pairs within the passel have the same user-supplied tag, sent message order type, transportation type, receive message order type, timestamp (if present), and set of sent region designators (if present). A passel is a message.}

\defentry{passive subscription}{A request to the runtime infrastructure (RTI) for the kinds of data (object classes and attributes as well as interactions) that the joined federate is currently interested in receiving. However, unlike an active subscription, this information is not used by the RTI to arrange for data to be delivered, nor is it used to tell publishing joined federates that another joined federate is subscribing to that data (by way of \hlaService{Start/Stop Registration}, \hlaService{Turn On/Off Updates}, or \hlaService{Turn On/Off Interactions} service invocations). This form of subscription is provided to support certain types of logger joined federates.}

\defentry{promoted}{Pertaining to an object instance, as known by a particular joined federate: having a discovered class that is a superclass of its registered class.}
\hlaNote{See 5.1.3 of IEEE Std 1516.1-2010.}

\defentry{publish}{To announce to a federation the information a federate may provide to the federation.}

\defentry{published}{}
\begin{hlalist}
    \item Pertaining to an object class from the perspective of a given joined federate: having at least one available attribute of the object class that was an argument, along with the object class, to a \hlaService{Publish Object Class Attributes} service invocation that was not subsequently unpublished via the \hlaService{Unpublish Object Class Attributes} service.
    \vspace{6.5pt}
    \hlaNote{See 5.1.2 of IEEE Std 1516.1-2010.}
    \vspace{3.5pt}
    \item Pertaining to an interaction class from the perspective of a given joined federate: being an argument to a \hlaService{Publish Interaction Class} service invocation that was not subsequently followed by an \hlaService{Unpublish Interaction Class} service invocation for that interaction class.
    \vspace{6.5pt}
    \hlaNote{See 5.1.3 of IEEE Std 1516.1-2010.}
\end{hlalist}

\newpage
\setlength{\parskip}{11.6pt}

\defentry{published attributes of an object class}{The class attributes that have been arguments to \textit{Publish Object}\linebreak
\textit{Class Attributes} service invocations by a given joined federate for that object class that have not\linebreak
subsequently been unpublished (either individually or by unpublishing the whole object class), and possibly\linebreak
the HLAprivilegeToDeleteObject attribute for that object class.}
\hlaNote{See 5.1.2 of IEEE Std 1516.1-2010.}

\defentry{range}{A subset of a dimension, defined by an ordered pair of values, the first being the range lower bound and the second being the range upper bound. This pair of values defines a semi-open interval [range lower bound, range upper bound) (i.e., the range lower bound is the smallest member of the interval, and the range upper bound is just greater than any member of the interval).}
\hlaNote{See 9.1.1 of IEEE Std 1516.1-2010.}

\defentry{range lower bound}{The first component of the ordered pair of values that is part of a range.}
\hlaNote{See 9.1 of IEEE Std 1516.1-2010.}

\defentry{range upper bound}{The second component of the ordered pair of values that is part of a range.}
\hlaNote{See 9.1 of IEEE Std 1516.1-2010.}

\defentry{received class}{The class that was an interaction's candidate received class at the joined federate when that interaction was received at that joined federate via an invocation of the \hlaService{Receive Interaction \textdagger}\ service.}
\hlaNote{See 5.1.3 of IEEE Std 1516.1-2010.}

\defentry{received parameters}{The set of parameters received when the \hlaService{Receive Interaction \textdagger}\ service is invoked. These parameters consist of the subset of the sent parameters of an interaction that are available parameters of the interaction's received class.}
\hlaNote{See 5.1.3 of IEEE Std 1516.1-2010.}

\defentry{receive order (RO)}{A characteristic of no ordering guarantee for messages. Messages that are received as RO messages will be received in an arbitrary order by the respective joined federate. A timestamp value will be provided with the message if one was specified when the message was sent, but that timestamp has no bearing on message receipt order.}

\defentry{reflect}{To receive new values for one or more instance attributes via invocation of the \hlaService{Reflect Attribute Values \textdagger}\ service.}
\hlaNote{See 6.7 of IEEE Std 1516.1-2010.}

\defentry{reflected instance attribute}{An instance attribute that is represented but not explicitly modeled in a joined federate. The reflecting joined federate accepts new values of the reflected instance attribute as they are produced by some other federation member and provided to it by the runtime infrastructure (RTI), via the \hlaService{Reflect Attribute Values \textdagger}\ service.}

\defentry{region}{A generic term that refers to either a region specification or a region realization. If not using data distribution management (DDM), region arguments may be ignored.}
\hlaNote{See 9.1 of IEEE Std 1516.1-2010.}

\defentry{region realization}{A region specification (set of ranges) that is associated with an instance attribute for update, with a sent interaction, or with a class attribute or interaction class for subscription. Region}

\newpage
\setlength{\parskip}{12pt}

realizations are created from region specifications via the \hlaService{Commit Region Modifications} (only when modifying a region specification from which region realizations are already derived), \hlaService{Register Object Instance With Regions}, \hlaService{Associate Regions for Updates}, \hlaService{Subscribe Object Class Attributes With Regions}, \hlaService{Subscribe Interaction Class With Regions}, \hlaService{Send Interaction With Regions}, or \hlaService{Request Attribute Value Update With Regions} services.
\hlaNote{See 9.1.1 of IEEE Std 1516.1-2010.}

\defentry{region specification}{A set of ranges. Region specifications are created using the \hlaService{Create Region} service, and a runtime infrastructure (RTI) is notified of changes to a region specification using the \hlaService{Commit Region Modifications} service.}
\hlaNote{See 9.1.1 of IEEE Std 1516.1-2010.}

\defentry{region template}{A region template is an incomplete region specification where one or more dimensions have not been assigned ranges.}

\defentry{register}{To invoke the \hlaService{Register Object Instance} or the \hlaService{Register Object Instance With Regions} service to create a unique object instance designator.}
\hlaNote{See 6.4 of IEEE Std 1516.1-2010.}

\defentry{registered class}{The object class that was an argument to the \hlaService{Register Object Instance} or the \hlaService{Register Object Instance With Regions} service invocation that resulted in the creation of the object instance designator for a given object instance.}

\defentry{regular interaction class description}{A description of an interaction class that follows IEEE Std 1516.2-\linebreak
2010 and contains at a minimum the interaction class name and publish/subscribe indicator.}

\defentry{regular object class description}{A description of an object class that follows IEEE Std 1516.2-2010 and contains at a minimum the object class name and publish/subscribe indicator.}

\defentry{remove}{To receive an invocation of the \hlaService{Remove Object Instance \textdagger}\ service for a particular object instance.}
\hlaNote{See 6.11 of IEEE Std 1516.1-2010.}

\defentry{repeated description}{A description of a construct (e.g., object class, interaction class, datatype, dimension) in a federation object model (FOM) module with a name and additional properties that are identical to a description that is already available in the current FOM.}

\defentry{resolution}{The smallest resolvable value separating attribute or parameter values that can be discriminated. Resolution may vary with magnitude for certain datatypes.}

\defentry{retraction}{An action performed by a federate to unschedule a previously scheduled message. Message retraction may be visible to the federate to whom the scheduled message was to be delivered. Retraction is widely used in classic event-oriented discrete event simulations to model behaviors such as preemption and interrupts.}

\defentry{runtime infrastructure (RTI)}{The software that provides common interface services during a High Level Architecture (HLA) federation execution for synchronization and data exchange.}

\defentry{runtime infrastructure (RTI) initialization data (RID)}{RTI vendor-specific information needed to run an RTI. If required, an RID is supplied when an RTI is initialized.}

\newpage
\setlength{\parskip}{12pt}

\defentry{scaffolding interaction class description}{A description of an interaction class in a federation object model (FOM) module or simulation object model (SOM) module that follows IEEE Std 1516.2-2010. However, it will contain only the name. The name of the scaffolding interaction class description is identical to the name of a regular interaction class description provided by another FOM/SOM module. Scaffolding interaction class descriptions act as placeholders in order to represent the class hierarchy from the other FOM/SOM module(s) from which a new regular interaction class description can be specified.}

\defentry{scaffolding object class description}{A description of an object class in a federation object model (FOM) module or simulation object model (SOM) module that follows IEEE Std 1516.2-2010. However, it will contain only the name. The name of the scaffolding object class description is identical to the name of a regular object class description provided by another FOM/SOM module. Scaffolding object class descriptions act as placeholders in order to represent the class hierarchy from the other FOM/SOM module(s) from which a new regular object class description can be specified.}

\defentry{sent class}{The interaction class that was an argument to the \hlaService{Send Interaction} or \hlaService{Send Interaction With Regions} service invocation that initiated the sending of a given interaction.}
\hlaNote{See 5.1.3 of IEEE Std 1516.1-2010.}

\defentry{sent interaction}{A specific interaction that is transmitted by a joined federate via the \hlaService{Send Interaction} or \hlaService{Send Interaction With Regions} service and received by other joined federates in the federation execution via the \hlaService{Receive Interaction \textdagger}\ service.}

\defentry{sent parameters}{The parameters that were arguments to the \hlaService{Send Interaction} or \hlaService{Send Interaction With Regions} service invocation for a given interaction.}
\hlaNote{See 5.1.3 of IEEE Std 1516.1-2010.}

\defentry{set}{A group of elements in which each element occurs at most once (this definition corresponds to the mathematical notion of sets). An example of a set would be a group of class attributes, each of which belongs to the same object class.}

\defentry{simulation object model (SOM)}{A specification of the types of information that an individual federate could provide to High Level Architecture (HLA) federations as well as the information that an individual federate can receive from other federates in HLA federations. The SOM is specified using one or more SOM modules. The standard format in which SOMs are expressed facilitates determination of the suitability of federates for participation in a federation.}

\defentry{simulation object model (SOM) module}{A subset of the SOM that contains some or all object model template (OMT) tables. SOM modules and one optional Management Object Model (MOM) and Initialization Module (MIM) are used to specify the SOM of a federate. A SOM module contains complete or scaffolding definitions for all object classes and interaction classes that are superclasses of object classes and interaction classes in the same SOM module.}

\defentry{specified dimensions}{The dimensions that are explicitly provided when the region specification is created or modified.}
\hlaNote{See 9.1.2 of IEEE Std 1516.1-2010.}

\defentry{stop publishing}{To take action that results in a class attribute that had been a published attribute of a class no longer being a published attribute of that class.}

\newpage
\setlength{\parskip}{11.4pt}

\defentry{subclass}{A class that adds additional detail to (specializes) another more generic class (superclass). A subclass, by inheriting the properties from its parent class (closest superclass), also inherits the properties of all superclasses of its parent as well.}

\defentry{subscribe}{To announce to a federation the information a federate wants from the federation.}

\defentry{subscribed}{}
\begin{hlalist}
    \item Pertaining to an object class from the perspective of a given joined federate: having subscribed attributes of that class or subscribed attributes of that class with regions, for some region. \hlaSeeAlso{\hlaTermRef{subscribed attributes of an object class}; \hlaTermRef{subscribed attributes of an object class with regions}.}
    \item Pertaining to an interaction class: being a subscribed interaction class or a subscribed interaction class with regions, for some region. \hlaSeeAlso{\hlaTermRef{subscribed interaction class}; \hlaTermRef{subscribed interaction class with regions}.}
\end{hlalist}

\defentry{subscribed attributes of an object class}{The class attributes that have been arguments to \hlaService{Subscribe Object Class Attributes} service invocations by a given joined federate for a given object class that have not subsequently been unsubscribed, either individually or by unsubscribing the whole object class.}
\hlaNote{See 5.1.2 and 5.6 of IEEE Std 1516.1-2010.}

\defentry{subscribed attributes of an object class with regions}{The class attributes that have been arguments to \hlaService{Subscribe Object Class Attributes With Regions} service invocations by a given joined federate for a given object class and a given region, assuming the joined federate did not subsequently invoke the \hlaService{Unsubscribe Object Class Attributes With Regions} service for that object class and region.}
\hlaNote{See 9.8 of IEEE Std 1516.1-2010.}

\defentry{subscribed interaction class}{An interaction class that, from the perspective of a given joined federate, was an argument to a \hlaService{Subscribe Interaction Class} or \hlaService{Subscribe Interaction Class With Regions} service invocation that was not subsequently followed by an \hlaService{Unsubscribe Interaction Class} or \hlaService{Unsubscribe Interaction Class With Regions} service invocation for that interaction class.}
\hlaNote{See 5.1.3 and 5.8 of IEEE Std 1516.1-2010.}

\defentry{subscribed interaction class with regions}{An interaction class and a region that, from the perspective of a given joined federate, were arguments to a \hlaService{Subscribe Interaction Class With Regions} service invocation that was not subsequently followed by an \hlaService{Unsubscribe Interaction Class With Regions} service invocation for that interaction class and that region.}
\hlaNote{See 9.10 of IEEE Std 1516.1-2010.}

\defentry{subscription region set}{A set of regions used for subscription of a class attribute or used for subscription of an interaction class. \hlaSeeAlso{\hlaTermRef{used for subscription of a class attribute}; \hlaTermRef{used for subscription of an interaction class}.}}

\defentry{superclass}{A class that generalizes a set of properties that may be inherited by more refined (i.e., detailed) versions of the class. In High Level Architecture (HLA) applications, a class may have at most one immediate superclass (i.e., can inherit only from a single class at the next higher level of the class hierarchy).}

\defentry{synchronization point}{A logical point in the sequence of a federation execution that all joined federates forming a synchronization set for that point attempt to reach and by which, if they are successful, they synchronize their respective processing.}

\newpage
\setlength{\parskip}{11.5pt}

\defentry{Time Advancing state}{The interval between a joined federate's request to advance its logical time and the corresponding grant by the runtime infrastructure (RTI). A joined federate may advance its logical time only by requesting a time advancement from the RTI via one of the following services:}
\begin{hlalist}
    \item \hlaService{Time Advance Request}
    \item \hlaService{Time Advance Request Available}
    \item \hlaService{Next Message Request}
    \item \hlaService{Next Message Request Available}
    \item \hlaService{Flush Queue Request}
\end{hlalist}

The joined federate's logical time will not actually be advanced until the RTI responds with a \hlaService{Time Advance Grant \textdagger}\ service invocation at that joined federate.

\defentry{time-constrained federate}{A joined federate that may receive timestamp order (TSO) messages and whose time advances are constrained by other joined federates within a federation execution.}
\hlaNote{See 8.1 of IEEE Std 1516.1-2010.}

\defentry{time management}{A collection of High Level Architecture (HLA) services that support controlled message ordering and delivery to the cooperating joined federates within a federation execution in a way that is consistent with federation requirements.}

\defentry{time-regulating federate}{A joined federate that may send timestamp order (TSO) messages and that constrains the time advances of other joined federates within a federation execution.}
\hlaNote{See 8.1 of IEEE Std 1516.1-2010.}

\defentry{timestamp (of message or save)}{The value of the timestamp argument provided to the relevant service invocation.}

\defentry{timestamp order (TSO)}{An ordering of messages provided by a runtime infrastructure (RTI) for joined federates making use of time management services and messages containing timestamps. Messages having different timestamps are said to be delivered in TSO if for any two messages, M1 and M2 (timestamped with T1 and T2, respectively), that are delivered to a single joined federate where T1 $<$ T2, then M1 is delivered before M2. Messages having the same timestamp will be delivered in an arbitrary order (i.e., no tie-breaking mechanism is provided by an RTI).}

\defentry{transportation type}{A runtime infrastructure (RTI) provided means of transmitting messages between joined federates. Different categories of service are defined with different characteristics such as reliability of delivery and message latency.}

\defentry{unspecified dimensions}{The available dimensions of a class attribute, instance attribute, interaction class, or sent interaction less the specified dimensions of the region specification from which the region realization is derived.}
\hlaNote{See 9.1.3 of IEEE Std 1516.1-2010.}

\defentry{update}{To invoke the \hlaService{Update Attribute Values} service for one or more instance attributes.}
\hlaNote{See 6.6 of IEEE Std 1516.1-2010.}

\defentry{update rate}{The rate at which instance attribute values are provided, either by an owning joined federate to the runtime infrastructure (RTI) or by the RTI to a subscribing joined federate.}

\newpage
\setlength{\parskip}{12pt}

\defentry{update region set}{A set of regions used for sending interactions or used for updating instance attributes.\linebreak \hlaSeeAlso{\hlaTermRef{used for sending}; \hlaTermRef{used for update}.}}

\defentry{used for sending}{}
\begin{hlalist}
    \item Pertaining to a region and the specified interaction class designator: being an argument in the \hlaService{Send Interaction With Regions} service.
    \item Pertaining to the default region: when the specified interaction class designator is an argument in the \hlaService{Send Interaction} service.
    \vspace{4pt}
    \hlaNote{See 9.1 of IEEE Std 1516.1-2010.}
\end{hlalist}
\vspace{-4pt}

\defentry{used for subscription of a class attribute}{}
\begin{hlalist}
    \item Pertaining to a region, an object class, and a class attribute: when the class attribute is a subscribed attribute of the object class with that region. \hlaSeeAlso{\hlaTermRef{subscribed attributes of an object class with regions}.}
    \vspace{6.5pt}
    \hlaNote{See 9.1 of IEEE Std 1516.1-2010.}
    \vspace{4pt}
    \item Pertaining to the default region: when the specified class attribute is a subscribed attribute of the specified class. \hlaSeeAlso{\hlaTermRef{subscribed attributes of an object class with regions}.}
    \vspace{6.5pt}
    \hlaNote{See 9.1 of IEEE Std 1516.1-2010.}
\end{hlalist}
\vspace{-4pt}

\defentry{used for subscription of an interaction class}{}
\begin{hlalist}
    \item Pertaining to a region and an interaction class: when the interaction class is a subscribed interaction class with regions. \hlaSeeAlso{\hlaTermRef{subscribed interaction class}.}
    \vspace{6.5pt}
    \hlaNote{See 9.1 of IEEE Std 1516.1-2010.}
    \vspace{4pt}
    \item Pertaining to the default region: when the specified interaction class is a subscribed interaction class. \hlaSeeAlso{\hlaTermRef{subscribed interaction class}.}
    \vspace{6.5pt}
    \hlaNote{See 9.1 of IEEE Std 1516.1-2010.}
\end{hlalist}
\vspace{-4pt}

\defentry{used for update}{}
\begin{hlalist}
    \item Pertaining to a region that, along with the specified object instance and attribute designators, has\linebreak been used as an argument in either the \hlaService{Register Object Instance With Regions} service or the\linebreak \hlaService{Associate Regions For Updates} service: when the region has not subsequently been used along with\linebreak the specified object instance designator as an argument in the \hlaService{Unassociate Regions For Updates}\linebreak service, nor has the joined federate subsequently lost ownership of the specified instance\linebreak attribute(s).
    \vspace{6.5pt}
    \hlaNote{See 9.1 of IEEE Std 1516.1-2010.}
    \vspace{4pt}
    \item Pertaining to the default region: when the specified instance attribute(s) is not currently used for update with any other region.
    \vspace{6.5pt}
    \hlaNote{See 9.1 of IEEE Std 1516.1-2010.}
\end{hlalist}
\vspace{-4pt}

\defentry{valid federate designator}{A federate designator which, during the federation execution, was returned by the \hlaService{Join Federation Execution} service to some joined federate. The federate does not have to remain a joined federate for its federate designator to remain a valid federate designator.}

\defentry{wall-clock time}{Time as determined by a chronometer such as a wristwatch or wall clock.}

\endgroup
